\section{Conclusão}

Neste laboratório implementamos em Java as transformações entre as coordenadas do mundo, as coordenadas normalizadas e as coordenadas do dispositivo, seguindo as fórmulas vistas em sala. O programa final acompanha a posição do mouse sobre o display o tempo todo, ao mover, clicar ou arrastar, e também permite digitar um ponto do mundo, escolher qual NDC usar e alterar a janela de mundo, mostrando em cada caso os valores de todos os sistemas envolvidos. Quando a coordenada sai dos limites do display, o programa mostra um alerta em vez de valores que não correspondem a nenhum pixel. Com isso, foi possível acompanhar na prática um caminho que, no quadro, parecia apenas uma sequência de regras de três.

Os testes mostraram que a implementação está de acordo com o que foi deduzido. O exemplo que resolvemos à mão na fundamentação, com a janela $[-2, 6] \times [-1, 4]$ e o ponto $(4, 2)$, chegou no programa ao mesmo pixel $(599, 359)$. Também conferimos que o NDC $[0,1]$ e o NDC $[-1,1]$ levam sempre ao mesmo pixel. No começo, esperávamos alguma diferença entre os dois, mas depois de montar as contas ficou claro que eles são só duas formas de escrever a mesma proporção. O mesmo vale para o arredondamento: ao voltar do pixel para o mundo, o ponto nunca é recuperado exatamente, e o erro fica sempre dentro de meio pixel, como era de se esperar.

Algumas questões só apareceram quando o programa começou a rodar. A principal foi a orientação do eixo $y$: no Java a origem fica no canto superior esquerdo e o $y$ cresce para baixo, enquanto em sala usamos a origem embaixo. Como não fizemos a inversão, a imagem no display fica espelhada na vertical. Outra foi a diferença de escala entre os eixos, já que a janela padrão é quadrada e o display não é. Nenhuma das duas altera as contas, mas elas mostram que, além das fórmulas, é preciso levar em conta como o dispositivo realmente organiza os pixels.

Como próximos passos, pensamos em inverter o eixo $y$ na conversão para o dispositivo, para que o display siga a mesma convenção da aula, e em manter a proporção entre os eixos ao definir a janela de mundo. Essas transformações também devem servir de base para os próximos trabalhos da disciplina, em que será necessário desenhar mais do que um único ponto na tela.
